\chapter{Título del segundo capítulo}%
\label{cap:capitulo2}%

\section{Contexto general de los peligros volcánicos}
El concepto de peligro volcánico comprende al conjunto de eventos que se producen en un volcán y que pueden provocar daños a personas o bienes expuestos (Ortiz y Araña, 1995), para cuya mitigación se realizan evaluaciones
y mapas para predecir las rutas que seguirán los eventos volcánicos. Entre los trabajos de referencia que tratan de la predicción y mitigación de los efectos de una erupción, destaca el de Fournier d’Albe (1979), 
quien define el peligro volcánico como la probabilidad de que un área determinada sea afectada por procesos o productos volcánicos potencialmente destructivos en un intervalo de tiempo dado.
\section{Lahares}
\subsection{Definición y mecanismos de formación}
La erupción paroxística de 2015 del Volcán Villarrica (Chile) produjo un lahar de 2,5 h de duración, que descendió más de 20 km dentro de la cuenca del Río Correntoso/Turbio y destruyó dos pequeños puentes. (Cita jonhson palma)
Los lahares son uno de los principales peligros en los volcanes. El lahar histórico más mortal ocurrió en 1985 en el Nevado del Ruiz (Colombia) cuando el derretimiento glacial inducido por la erupción produjo flujos 
que viajaron más de 100 km y mataron a más de 25,000 personas \cite{Lowe1986}.
\subsection{Clasificación de los lahares}
\subsection{Comportamiento dinámico y reológico} %No se si deberia mencionar la reologia o no

\section{Volcán Villarrica}
\subsection{Contexto geográfico y geológico}
El volcán Villarrica corresponde a uno de los más de 200 volcanes pleistocenos a holocenos ubicados en el arco Andino, los cuales están agrupados en los siguientes segmentos:
Zona Volcánica Norte (2°N-5°N), Zona Volcánica Central (14°S-28°S), Zona Volcánica Sur (33°S-46°S) y Zona Volcánica Austral (49°S-55°S) \cite{article}.
\subsection{Actividad histórica reciente}
\subsection{Justificación como caso de estudio}

\section{Riesgo y exposición}
\subsection{Zonas de amenaza}
\subsection{Antecedentes de gestión del riesgo}

\section{Modelamiento computacional de lahares (estado del arte)}
\subsection{Enfoques tradicionales}
\subsection{Brecha que aborda esta tesis}

Utilice \lstinline!\chapter! para agregar nuevos capítulos, de este modo se numerarán en forma correlativa y la tabla de contenido los capturará correctamente.

Utilice tablas para sintetizar datos. Estas deberían quedar referenciadas en un índice en las páginas iniciales, al igual que las figuras, al colocarles una etiqueta con \lstinline!\label!. \parencite{MR1487846}\par%

\medskip%

\begin{table}[!h]
	\Centering
	\caption{Resultados partículas exóticas}%
	\label{tab:tabla2-1}%
	\begin{tabular}{ll}
		\textbf{Parámetro del experimento} & \textbf{Valor obtenido} \\ \hline
		Energía de colisión                & 14 TeV                  \\ \hline
		Partículas detectadas              & Masa de 200 GeV/c²      \\ \hline
		Sección transversal                & 0.01 pb                
	\end{tabular}
\end{table}%

Entregue independencia a párrafos que contienen teoremas o fórmulas matemáticas, insertando las ecuaciones de este modo:\par%

\vspace{10mm}
\begin{theorem}[Convergencia Monótona]
\label{thm:conv-mon}
	Sea $\{f_n\}$ una sucesión de funciones medibles tales que:
	1. $f_n(x) \geq 0$ para todo $x$ y todo $n$.
	2. $f_n(x)$ es monótona creciente, es decir, $f_n(x) \leq f_{n+1}(x)$.
	3. Existe una función límite $f(x) = \lim_{n \to \infty} f_n(x)$.
	
	Entonces, se cumple que:
	\begin{equation}
		\lim_{n \to \infty} \int f_n(x) \,dx = \int \lim_{n \to \infty} f_n(x) \,dx.
	\end{equation}
\end{theorem}%

\begin{proof}
	Dado que $\{f_n\}$ es creciente, tenemos:
	\begin{equation}
	    f_1(x) \leq f_2(x) \leq \dots \leq f(x).
	\end{equation}

	Por la propiedad de la integral, la sucesión $I_n = \int f_n(x) \,dx$ es monótona creciente y acotada, por lo que tiene un límite finito:
	\begin{equation}
	    \lim_{n \to \infty} I_n = I, \quad \text{donde } I = \int f(x) \,dx.
	\end{equation}

Aplicando la convergencia puntual de funciones medibles:
\begin{equation}
    \lim_{n \to \infty} \int f_n(x) \,dx = \int \lim_{n \to \infty} f_n(x) \,dx.
\end{equation}

Esto completa la demostración. \qed
\end{proof}%
\vspace{10mm}

Agregue fragmentos de código con \lstinline!\begin{lstlisting}!.

\begin{lstlisting}[language=Python,caption={Hello, World!},label={lst:hello}]
def main():
	print("Hello, World!")

if __name__ == "__main__":
	main()
\end{lstlisting}

Otra cita bibliográfica interesante, a modo de ejemplo. \parencite{MR1823409}%